\subsection{Significance of the Photodiode: Sampling Rate, Velocity Range, and System Precision}
\label{sec:importance_bpw34_precision}

\subsubsection{Maximum Measurable Velocity with Camera}
Because the camera samples the irradiance waveform at $f_s = 30\text{ Hz}$, fringe skipping (temporal aliasing) occurs when the fringe-crossing rate $f_{\text{fringe}}$ becomes too fast to resolve between consecutive frames:
\begin{itemize}
    \item \textbf{Nyquist Limit for Directional Phase Tracking:} To avoid aliasing of a sinusoidally modulated signal and preserve direction reconstruction from frame samples, the fringe frequency must satisfy $f_{\text{fringe}} < \frac{f_s}{2} = 15\text{ Hz}$ (at most $\frac{1}{2}$ fringe cycle per frame). This sets the maximum velocity at
    \begin{equation}
        v_{M2, \max}^{\text{(C)}} = \frac{\lambda}{2}\,\frac{f_s}{2} = \frac{\lambda f_s}{4}
        \label{eq:vmax_nyquist}
    \end{equation}
\end{itemize}

The setup laser source has a  $\lambda = 632.8\text{ nm}$ and with the camera framrate we obtain:

\begin{equation}
    v_{M2, \max}^{\text{(C)}} \approx \frac{632.8 \times 10^{-9}\text{ m} \times 30\text{ s}^{-1}}{4} \approx 4.75 \times 10^{-6}\text{ m/s} = 4.75\ \mu\text{m/s}
\end{equation}


If mirror $M2$ moves faster than these limits, fringes cross the $\text{ROI}$ faster than the $30\text{ fps}$ frame rate can sample them, causing missed maxima/minima, underestimated displacement, and direction reversal artifacts (aliasing). Motion blur from the camera exposure time also reduces fringe contrast at higher speeds.

\subsubsection{Maximum Measurable Velocity Photodiode}
\label{sec:vmax_bpw34_fpga}

Using the PIN photodiode BPW34 (\cref{sec:PIN_Photodiode}) operated in reverse-biased mode, temporal sampling is no longer limited by a $30\text{ Hz}$ camera frame rate. Instead, fringe skipping occurs when the fringe-crossing frequency $f_{\text{fringe}}$ exceeds the effective analog bandwidth governed by the BPW34 carrier response, the TIA feedback time constant, the comparator transition, and the FPGA input-pin load, detailed in \cref{sec:circuit_schematic}.

Because the comparator outputs a single binary logic stream, direction ambiguity is not resolved from one channel alone; thus, the measurement yields unsigned fringe-crossing counts ($\Delta m \ge 0$, i.e., speed $|v_{M2, \max}|$, or signed velocity when direction is externally constrained).

A full fringe cycle requires at least one rising and one falling transition ($\sim 2 t_{\text{sw}}$), while sinusoidal attenuation in the TIA limits peak-to-peak swing above $f_{-3\text{dB}}$. Hence, the maximum resolvable fringe-crossing frequency is bounded by
\begin{equation}
    f_{\text{fringe, max}}^{(\text{PD})} \approx \min\!\left(f_{-3\text{dB}},\, \frac{1}{2 t_{\text{sw}}}\right)
    \label{eq:ffringe_max_bpw34}
\end{equation}
and, since each fringe cycle corresponds to an axial displacement of $\frac{\lambda}{2}$ for movable mirror $M2$, the maximum measurable mirror speed before fringe skipping is
\begin{equation}
    |v_{M2, \max}^{(\text{PD})}| = \frac{\lambda}{2}\,f_{\text{fringe, max}}^{(\text{PD})}
    \label{eq:vmax_bpw34_fpga}
\end{equation}

We have $f_{-3\text{dB}} \approx 6.37\text{ MHz}$, reducing its first-order $10\%\text{--}90\%$ rise/fall time to
\begin{equation}
    t_{r,\text{TIA}} = 2.2\,R_f C_f = 2.2 \times (10\,\text{k}\Omega)(2.5\text{ pF}) = 55\text{ ns}
\end{equation}
The remaining signal-path stages contribute:
\begin{itemize}
    \item \textbf{BPW34 carrier response:} $t_{r,\text{PD}} \approx 100\text{ ns}$ (representative datasheet transition time, now the dominant speed bottleneck).
    \item \textbf{TLV3542 comparator:} across a $3.3\text{ V}$ swing at $150\text{ V}/\mu\text{s}$ we obtain: \\
    $t_{\text{slew}} \approx 3.3\text{ V} / (150\text{ V}/\mu\text{s}) \approx 22\text{ ns}$.
    \item \textbf{Arty Z7-20 I/O pin load:} $R_{\text{out}} = 47\,\Omega$ in series with the board's protection resistor $200\,\Omega$ on Shield digital I/O pins, driving $C_{\text{load}} \approx 20\text{ pF}$ yields at most $t_{r,\text{pin}} \approx 2.2\,(247\,\Omega)(20\text{ pF}) \approx 10.9\text{ ns}$.
\end{itemize}
Combining all four stages by root-sum-of-squares gives the total front-end switching time:
\begin{equation}
    t_{\text{sw}} \approx \sqrt{t_{r,\text{TIA}}^2 + t_{r,\text{PD}}^2 + t_{\text{slew}}^2 + t_{r,\text{pin}}^2}
    \approx \sqrt{55^2 + 100^2 + 22^2 + 10.9^2}\text{ ns} \approx 117\text{ ns}
\end{equation}
Because the TIA bandwidth ($6.37\text{ MHz}$) now exceeds the switching-time limit, the maximum resolvable fringe frequency is governed by $\frac{1}{2 t_{\text{sw}}}$:
\begin{equation}
    f_{\text{fringe, max}}^{(\text{PD})} \approx \frac{1}{2 \times 117\text{ ns}} \approx 4.27\text{ MHz}
\end{equation}
For a red laser at $\lambda = 632.8\text{ nm}$, the corresponding maximum measurable mirror speed is:
\begin{equation}
    |v_{M2, \max}^{(\text{PD})}| \approx \frac{632.8 \times 10^{-9}\text{ m}}{2 \times 1} \times 4.27 \times 10^6\text{ Hz} \approx 1.35\text{ m/s}
\end{equation}

\subsubsection{System Precision}

Integrating the reverse-biased BPW34 photodiode and the FPGA counter transforms the interferometer from a quasi-static imaging setup into a high-speed metrological instrument. While the camera captures two-dimensional interference patterns, its $33.3\text{ ms}$ frame interval introduces motion blur and limits fringe tracking speed. By collapsing detection to a fast single-point photodetector sampled continuously in programmable logic on the FPGA, the photodiode setup overcomes this temporal bottleneck in three critical ways:

\begin{itemize}
    \item \textbf{Higher Effective Sampling Rate:}
    Instead of integrating photons over tens of milliseconds, the photodiode setup converts optical irradiance variations into real-time digital voltage swings with a combined rise time of $t_{\text{sw}} \approx 117\text{ ns}$.

    \item \textbf{Velocity Measurement:}
    Raising the maximum fringe-crossing frequency from $15\text{ Hz}$ to $4.27\text{ MHz}$ extends the measurable mirror speed for $\lambda = 632.8\text{ nm}$ from $4.75\ \mu\text{m/s}$ up to $|v_{M2, \max}^{(\text{PD})}| \approx 1.35\text{ m/s}$, an expansion of nearly $3 \times 10^5$. This dynamic range allows the interferometer to track motorized linear actuators, rapid transient displacements, and high-frequency mechanical vibrations without fringe skipping. Moreover, because the FPGA can latch a free-running $8\text{ ns}$ timestamp counter on each incoming edge, instantaneous mirror speed can be reconstructed on a per-fringe basis via $|v_{M2}(t)| = \frac{\lambda}{2\,\Delta t}$, where $\Delta t$ is the measured period between consecutive fringe transitions.

    \item \textbf{Sub-Micrometer Spatial and Nanosecond Temporal Precision:}
    The overall precision of the system arises from the combination of optical interferometric resolution, clean analog conditioning, and synchronous digital capture:
    \begin{enumerate}
        \item \textit{Optical quantization:} Each complete fringe cycle corresponds to a mirror displacement of $\frac{\lambda}{2} = 316.4\text{ nm}$ (or $\frac{\lambda}{4} = 158.2\text{ nm}$ when the FPGA edge detector counts both rising and falling transitions), providing intrinsic nanometer-scale displacement without requiring analog-to-digital amplitude interpolation.
        \item \textit{Temporal quantization:} At a $125\text{ MHz}$ FPGA clock, timing jitter due to digital sampling is negligible across typical mechanical speeds. For example, at a mirror velocity of $10\text{ mm/s}$ ($f_{\text{fringe}} \approx 31.6\text{ kHz}$, $T_{\text{fringe}} \approx 31.6\,\mu\text{s}$), the $8\text{ ns}$ clock resolution corresponds to a relative timing precision of roughly $0.025\%$ per fringe period.
        \item \textit{Noise Reduction:} Operating the TIA just above critical damping ($C_f = 2.5\text{ pF} > C_{f,\min}$) avoids ringing on fast optical transients, while the $150\text{ mV}$ Schmitt-trigger hysteresis band avoids unwanted digital changes when crossing the voltage thresholds.
    \end{enumerate}
\end{itemize}

In summary, whereas a camera excels at capturing two-dimensional spatial fringe geometry for static wavefront inspection, the \textbf{BPW34} + \textbf{TLV3542} + \textbf{Arty Z7-20} setup provides the high temporal bandwidth, wide velocity ceiling, and noise-resilient nanosecond precision needed for real-time dynamic displacement and speed metrology (with direction either externally constrained or inferred from known actuator trajectories).